% ---------------------------------------------------------------------
%  Sample content: text, lists and mathematics
% ---------------------------------------------------------------------
\section{Text and lists}\label{sec:text}

This document is a sample built with the template. Each section shows one
kind of element, so you can see the final look before writing your own notes.
Replace this text with your content.

Body text supports \emph{emphasis}, \textbf{bold text},
\hl{highlighted key terms}, inline code such as \code{result = compute(data)},
and links such as \href{https://example.com}{example.com}.\footnote{Footnotes
appear at the bottom of the page.}

\begin{itemize}
  \item Bullet lists use the accent color.
  \item They can be nested:
  \begin{itemize}
    \item a second-level item
    \item another second-level item
  \end{itemize}
  \item[\textcolor{co-green}{\checkmark}] Custom markers are possible too.
\end{itemize}

\begin{enumerate}
  \item Numbered lists
  \item keep the spacing compact
  \item and align neatly.
\end{enumerate}

\subsection{A subsection}

Subsections and subsubsections share the accent color, with smaller sizes.
The header of each page shows the current section.

\subsubsection{A subsubsection}

Paragraphs are separated by a small gap instead of an indent, which keeps
notes easy to scan.

% ---------------------------------------------------------------------
\section{Mathematics}\label{sec:math}

Inline math sits in the text: $f\colon \R^n \to \R$, $\norm{\vect{x}}$,
$\abs{a - b}$ and $\inner{\vect{u}}{\vect{v}}$. Letters and digits use the same
typeface as the text.

A numbered equation can be referenced later:
\begin{equation}\label{eq:sample}
  \int_a^b f(x)\,dx = F(b) - F(a).
\end{equation}

A multi-line derivation, aligned at the equals sign, with only the last line
numbered:
\begin{align}
  (a + b)^2 &= (a + b)(a + b) \nonumber\\
            &= a^2 + 2ab + b^2. \label{eq:aligned}
\end{align}

Matrices, sums, fractions and cases:
\[
  A = \begin{bmatrix} a_{11} & a_{12} \\ a_{21} & a_{22} \end{bmatrix},
  \qquad
  \sum_{k=1}^{n} k = \frac{n(n+1)}{2},
  \qquad
  \abs{x} = \begin{cases} x, & x \ge 0, \\ -x, & x < 0. \end{cases}
\]
References such as \eqref{eq:sample} and \eqref{eq:aligned} are clickable links.
